% Candidate theorem section; no submission or novelty claim.
\section{Scale-invariant local-price search}
Let $C$ be the shared boundary of an explicit e-graph, $s=|C|$, and let
$c_v(S)$ be the inherited minimum local construction price. Costs are
nonnegative integers and the threshold is $K$. Define $P_v$ to contain every
price at most $K$ of a private local implementation of $v$, stopping at
boundary references. Boundary references are not assumed constructed.
Let $A_v=\{0\}\cup P_v$.

\paragraph{Snapping lemma.}
For $0\leq b_v\leq K$, define
\[
 \pi_v(b_v)=\max\{a\in A_v:a\leq b_v\}.
\]
Then, for every $S\subseteq C$,
\[
 c_v(S)\leq b_v\quad\Longleftrightarrow\quad c_v(S)\leq\pi_v(b_v).
\]
Indeed every finite minimum below $K$ is realized by a local implementation
and hence belongs to $P_v$; the reverse direction follows from
$\pi_v(b_v)\leq b_v$. Grounding from the empty set therefore has exactly
the same synchronous rounds under $b$ and $\pi(b)$, with no increase in total
budget. The grounded-budget threshold theorem remains exact after restricting
to the finite alphabets. The alphabets need not be minimal: an attainable
local price may never be optimal under any boundary availability pattern.

\paragraph{Mandatory grounding.}
Let $M\subseteq C$ be any sound set of mandatory classes. Set
$\ell_v=c_v(C\setminus\{v\})$ for $v\in M$ and $\ell_v=0$ otherwise.
Every successful budget vector satisfies $b\geq\ell$.
Let $G$ be the least grounded set under $\ell$. Replacing $b_v$ by $\ell_v$
for $v\in G$ preserves its final grounded set and cannot raise cost.
To see this, the new budgets still dominate $\ell$, so $G$ remains available;
all other enabling tests retain their caps. Induction through the old
construction proves equality of the final sets. This justifies fixing those
coordinates before price compilation. It does not assert equality of every
intermediate round after this second reduction.

\paragraph{Quantum corollary with charged preprocessing.}
Suppose the surviving alphabets have product $Q$ and total explicit table size
$P$, and their computation costs $\mathcal A$. For $Q\geq2$, a binary index
of length $\lceil\log_2 Q\rceil$ is reversibly decoded by mixed-radix division
and explicit table scans. Reject padding, check the sum, perform grounded
construction, mark, and uncompute. Standard bounded-error quantum search gives
\[
 \mathcal A+O\!\left(\sqrt Q\,\operatorname{poly}(N,P,\log(K+2))\right)
\]
time and polynomial workspace in the compiled representation. Zero/one-label
domains are settled classically. No QRAM or exponential lookup table is assumed
free. A measured index produces an independently checked classical program.
Multiplying every cost and $K$ by an integer $\lambda>0$ multiplies the prices,
not $Q$. Arithmetic bit costs still grow with $\log\lambda$.

\paragraph{Compilation obstruction.}
For $i=0,\ldots,n-1$, let a root require private $p_i,q_i$, each obtainable
for price $3^i$ or for zero using a shared $b_i$ of price $2\cdot3^i+1$.
The root is zero-cost. Its local price set is
$\{0,\ldots,3^n-1\}$; its conditional minimum prices number $2^n$.
No class grounds at the zero lower vector. Yet the exact optimum is
$\sum_i2\cdot3^i=3^n-1$, because independent direct construction is cheaper
than sharing in every component. Thus even distinct local thresholds can be
exponentially numerous on a classically easy graph. Explicit preprocessing
must be charged, and this corollary does not promise polynomial $\mathcal A$
or $P$ in the original input.

\paragraph{Comparison boundary.}
The $5^w$ treewidth bound of Goharshady--Lam--Parreaux is not, unchanged,
a bound for acyclic extraction from cyclic alternatives; their discussion
states the reachability extension incurs a $c^{w^2}$ dependence. The
acyclicity-aware construction of Sun--Zhang--Ni is a relevant
$2^{O(w^2)}\operatorname{poly}(N,w)$ comparator. Neither this distinction
nor the present corollary establishes superiority over those methods, the
$2^s$ subset dynamic program, representative enumeration, or native solvers.
